are simplified considerably by making a particular choice of ``gauge''---the Regge-Wheeler
choice---for the gravitational field. Making such a choice of gauge corresponds to removing
the coordinate arbitrariness in $(t, r, \theta, \phi)$, which arbitrariness is of the same
size as the perturbation, $h_{\mu\nu}$, in the geometry of spacetime.

The equations of motion can be simplified still further by restricting one's attention
to spherical harmonics with $M = 0$. Once these are understood, the spherical-harmonic
motions ($l$, $M \neq 0$, $r$) can be obtained from $(l$, $M = 0$, $r)$ by suitable rotations about
the center of the star.

Having specialized to a particular spherical harmonic ($l$, $M$, $r$), having made the
Regge-Wheeler choice of gauge for that spherical harmonic, and having restricted attention
to $M = 0$, we obtain the following vastly simplified forms for the fluid displacement,
$\xi$, and the metric perturbation, $h_{\mu\nu}$ (see Appendix~A for derivation):

If $r = (-1)^{l+1}$---``odd parity''; ``magnetic-type'' perturbations---then the covariant
components of the fluid-displacement vector take the form
\begin{equation}
  \xi_{t} = \xi_{r} = 0 \;, \qquad
  \xi_{\theta} = U(r,t)\,\sin\theta\,\partial_{\phi} P_{l}(\cos\theta) \;;
\label{eq:6a}\tag{6a}
\end{equation}
and the only non-vanishing components of the metric perturbation are
\begin{align}
  h_{t\phi} &= h_{\phi t} = h_{0}(r,t)\,\sin\theta\,\partial_{\theta} P_{l}(\cos\theta) \;,\notag\\
  h_{r\phi} &= h_{\phi r} = h_{1}(r,t)\,\sin\theta\,\partial_{\theta} P_{l}(\cos\theta) \;.
\label{eq:6b}\tag{6b}
\end{align}

If $r = (-1)^{l}$---``even parity''; ``electric-type'' perturbations---then the contravariant
components of the fluid displacement vector take the form
\begin{equation}
  \xi^{r} = r^{-2} e^{-\lambda/2} V(r,t)\, P_{l}(\cos\theta) \;, \qquad
  \xi^{\theta} = -r^{-2} W(r,t)\,\partial_{\theta} P_{l}(\cos\theta) \;, \qquad
  \xi^{\phi} = 0 \;;
\label{eq:7a}\tag{7a}
\end{equation}
and the metric perturbation takes the form
\begin{equation}
  h_{\mu\nu} =
  \left(
  \begin{array}{c|cccc}
     & t & r & \theta & \phi \\
    \hline
    t & H_{0} e^{\nu} & H_{1} & 0 & 0 \\
    r & H_{1} & H_{2} e^{\lambda} & 0 & 0 \\
    \theta & 0 & 0 & r^{2} K & 0 \\
    \phi & 0 & 0 & 0 & r^{2} \sin^{2}\!\theta\, K
  \end{array}
  \right)
  P_{l}(\cos\theta) \;.
\label{eq:7b}\tag{7b}
\end{equation}

Here $V$, $W$, $H_{0}$, $H_{1}$, $H_{2}$, and $K$ are functions of $r$ and $t$.

For odd-parity perturbations the equations of motion are a set of coupled equations for
the fluid-displacement function $U(r,t)$ and the metric perturbation functions $h_{0}(r,t)$,
$h_{1}(r,t)$ of equations~(6). For even-parity perturbations the equations of motion are a set
of coupled equations for the fluid-displacement functions $V(r,t)$, $W(r,t)$ and the metric
perturbation functions $H_{0}(r,t)$, $H_{1}(r,t)$, $H_{2}(r,t)$, $K(r,t)$ of equations~(7). The odd and even
cases are considered separately in the next two sections.

\subsection*{(c) Odd-Parity Motions}

In Appendix~B the Einstein field equations are calculated and discussed for the odd-parity
case. From that discussion one learns that the odd-parity motions of a star are not
characterized by pulsations which emit gravitational waves; rather, they are characterized
by a stationary, differential rotation of the fluid inside the star and by gravitational
waves which do not couple to the star at all.

This result should not be surprising. Pulsations can occur only if the perturbation
causes a change in the star's internal density and pressure, $\rho$ and $p$. However, $\rho$ and $p$
are scalar fields; and all \textit{scalar} spherical harmonics are of even parity. (Only vector and
tensor spherical harmonics can have odd parity.) Therefore, a perturbation of odd
tensor spherical harmonics cannot change the star's density or pressure distributions and therefore cannot
cause the star to pulsate.
